\numberwithin{equation}{section}

\newtheorem{thm}{Theorem}[section]
\newtheorem*{Theorem*}{Theorem}
\newtheorem{cor}[thm]{Corollary}
\newtheorem{lem}[thm]{Lemma}
\newtheorem{prop}[thm]{Proposition}

\newtheorem{theoremalpha}{Theorem}
\renewcommand*{\thetheoremalpha}{\Alph{theoremalpha}} % Theorem in the Intro
\newtheorem{corollaryalpha}[theoremalpha]{Corollary} % Corollary in the Intro
\newtheorem{questionalpha}[theoremalpha]{Question} % Question in the Intro
\newtheorem{propositionalpha}[theoremalpha]{Proposition} % Proposition in the intro

\theoremstyle{definition}
\newtheorem{Definition}[thm]{Definition}
\newtheorem{Note}[thm]{Note}
\newtheorem{Example}[thm]{Example}
\newtheorem{rem}[thm]{Remark}

\usepackage[all]{xy} % Xy-Pic for diagrams (disabled)
\usepackage{tikz}
\usetikzlibrary{cd}
\usetikzlibrary{decorations.pathmorphing}

\usepackage{enumitem}


\usepackage[normalem]{ulem}

\DeclareMathOperator{\supp}{supp}
\DeclareMathOperator{\curv}{curv}
\DeclareMathOperator{\diag}{diag}
\DeclareMathOperator{\diam}{diam}
\DeclareMathOperator{\Dif}{Diff}
\DeclareMathOperator{\Fix}{Fix}
\DeclareMathOperator{\Iso}{Isom}
\DeclareMathOperator{\dist}{dist}
\DeclareMathOperator{\inte}{int}
\DeclareMathOperator{\img}{Image}
\DeclareMathOperator{\codim}{codim}
\DeclareMathOperator{\frk}{Erank}
\DeclareMathOperator{\Tub}{Tub}
\DeclareMathOperator{\rank}{rank}
\DeclareMathOperator{\rk}{rk}
\DeclareMathOperator{\Ker}{Ker}
\DeclareMathOperator{\Susp}{Susp}
\DeclareMathOperator{\Span}{span}
\DeclareMathOperator{\Ric}{\mathrm{Ric}}
\DeclareMathOperator{\Home}{\mathrm{Homeo}}
\DeclareMathOperator{\maxsymrk}{maxsymrk}
\DeclareMathOperator{\vol}{vol}
\DeclareMathOperator{\rad}{rad}
\DeclareMathOperator{\Pol}{Pol}
%\DeclareMathOperator{\prin}{prin}

\usepackage{tikz}

\newcommand{\N}{\mathbb{N}} %Naturals
\newcommand{\R}{\mathbb{R}} %Reals
\newcommand{\Z}{\mathbb{Z}} %Integers
\newcommand{\C}{\mathbb{C}} %Complex
\newcommand{\D}{\mathbb{D}}
\newcommand{\Sp}{S} %Sphere
\newcommand{\RP}{\mathbb{R}P} % Real projective space
\newcommand{\CP}{\mathbb{C}P} %Complex Projective space
\newcommand{\Id}{\mathrm{Id}}
\newcommand{\Co}{C} %Differentiable in degree n
\newcommand{\symrk}{\mathrm{symrk}}
\newcommand{\Tau}{\mathrm{T}}
\newcommand{\U}{\mathrm{U}}
\newcommand{\prin}{\mathrm{prin}}

\newcommand{\fol}{\mathcal{F}}
\newcommand{\blue}[1]{{\color{blue} #1}}

\newcommand{\p}{\pi}
\newcommand{\Rho}{\mathrm{P}}

% Custom function restriction command % \restr{f}{A}
% See https://tex.stackexchange.com/questions/22252/how-to-typeset-function-restrictions
\newcommand\restr[2]{{% we make the whole thing an ordinary symbol
		\left.\kern-\nulldelimiterspace % automatically resize the bar with \right
		#1 % the function
%		\vphantom{\big|} % pretend it's a little taller at normal size
		\right|_{#2} % this is the delimiter
}}


\allowdisplaybreaks
\usepackage[a4paper,margin=20mm]{geometry}
\AddToHook{cmd/maketitle/after}{\markboth{Corro and Galaz-Garcia}{Foliated isometry dimension and rigidity}}
\pdfstringdefDisableCommands{\def\boldsymbol#1{#1}\def\bm#1{#1}\def\mathbb#1{#1}}
